\subsection{导子与同态的提升}

设 \(k\) 为环，C 为 \(k\)-代数，N 为 C 的平方零理想。以 \(\pi : \mathrm{C} \to \mathrm{C/N}\) 表示典范同态；由于 \(\mathrm{N}^2 = \{0\}\)，N 的 C-模结构来自一个 C/N-模结构。

设 A 为 \(k\)-代数，\(\varphi : \mathrm{A} \to \mathrm{C/N}\) 为 \(k\)-代数同态。赋予 N 由 \(\varphi\) 导出的 A-模结构。我们称一个 \(lifting\)（提升；到 C 中）指的是 \(\varphi\) 的任一 \(k\)-代数同态 \(\tilde{\varphi} : \mathrm{A} \to \mathrm{C}\)，它满足 \(\pi \circ \tilde{\varphi} = \varphi\)。设 \(\tilde{\varphi}\) 为这样一个提升，\(\delta\) 为从 A 到 N 中的映射；以 \(\delta + \tilde{\varphi}\) 表示映射 \(x \mapsto \delta(x) + \tilde{\varphi}(x)\)（从 \(A\) 到 C 中）。

命题 1.--- 若 φ 有一个提升，则映射 (δ, φ) ↦ δ + φ 定义了 A 到 N 中的 k-导子群在 φ 的提升集合上的一个单可递作用。

设 \(\tilde{\varphi}_{0}: A \to C\) 为 \(\varphi\) 的一个提升。映射 \(\delta \mapsto \delta + \tilde{\varphi}_{0}\) 诱导出从 \(A\) 到 \(N\) 中的映射的集合到映射 \(\tilde{\varphi}: A \to C\)（满足 \(\pi \circ \tilde{\varphi} = \varphi\)）组成的集合上的一个双射。固定 \(\delta\)，令 \(\tilde{\varphi} = \delta + \tilde{\varphi}_{0}\)。为使 \(\tilde{\varphi}\) 是 \(k\) -代数同态，必须且只需 \(\delta\) 是一个 \(k\)-导子（自 \(A\) 到 \(N\) 中）：事实上，对 \(x, y\)（\(A\) 中的）与 \(\lambda\)（\(k\) 中的），有关系式

\[
\begin{array}{c} \tilde {\varphi} (x + y) - \tilde {\varphi} (x) - \tilde {\varphi} (y) = \delta (x + y) - \delta (x) - \delta (y) \\ \tilde {\varphi} (\lambda x) - \lambda \tilde {\varphi} (x) = \delta (\lambda x) - \lambda \delta (x) \\ \tilde {\varphi} (x y) - \tilde {\varphi} (x) \tilde {\varphi} (y) = \delta (x y) - \delta (x) \delta (y) - \delta (x) \tilde {\varphi} _ {0} (y) - \tilde {\varphi} _ {0} (x) \delta (y) \\ = \delta (x y) - \varphi (x) \delta (y) - \varphi (y) \delta (x), \end{array}
\]

最后一处等式由 N 的平方为零这一事实得出。命题得证。

例.--- 设 B 为 k-代数，N 为 B-模。赋予 k-模 B⊕N 以由 \((b,x)(b',x') = (bb',bx' + b'x)\) 定义的 k-代数结构 (参见 A, III, 第 127 页)，于是 N 是 B ⊕ N 中的平方零理想。设 \(\varphi : A \to B\) 为 k-代数同态。于是 \(\varphi\) 到 B ⊕ N 中的提升恰是形如 \(x \mapsto (\varphi(x), \delta(x))\) 的映射，其中 \(\delta\) 取遍 A 到 N 中的 k-导子的集合 (参见同上引文, 命题 12)。

设 \(\Omega_{k}(A)\) 为环 A 的 \(k\)-微分模，\(d: \mathrm{A} \longrightarrow \Omega_{k}(\mathrm{A})\) 为泛 \(k\)-导子 (A, III, 第 133–134 页)；回忆 (同上引文)：对每个 A-模 M，映射 \(v \mapsto v \circ d\) 是从 \(\mathrm{Hom}_{\mathrm{A}}(\Omega_{k}(A), \mathrm{M})\) 到 A 到 M 中的 \(k\)-导子组成的 A-模上的 A-线性同构。

设 J 为 A 的理想。据 A, III, 第 137 页，有 A/J-线性映射的正合列

\[
\mathrm{J} / \mathrm{J} ^ {2} \stackrel {\bar {d}} {\longrightarrow} (\mathrm{A} / \mathrm{J}) \otimes_ {\mathrm{A}} \Omega_ {k} (\mathrm{A}) \longrightarrow \Omega_ {k} (\mathrm{A} / \mathrm{J}) \rightarrow 0,
\]

其中 \(\bar{d}\) 是由 \(d\) 在 J 上的限制过渡到商而导出的同态。

设 \(\rho : A \to A/J^{2}\) 与 \(\pi : A/J^{2} \to A/J\) 为典范满射。设 \(v : (A/J) \otimes_{A} \Omega_{k}(A) \longrightarrow J/J^{2}\) 为一个 \(k\)-线性映射；对它结合一个 \(k\)-线性映射 \(H_{v} : A \to A/J^{2}\)，令 \(H_{v}(x) = \rho(x) - v(1 \otimes dx)\)。若 \(v\) 是 \(\bar{d}\) 的一个收缩，则 \(H_{v}\) 在 \(J\) 上为零，从而通过过渡到商诱导出 \(k\)-线性映射 \(h_{v} : A/J \to A/J^{2}\)。另一方面，给定 \(k\)-线性映射 \(h : A/J \to A/J^{2}\)，以 \(\psi_{h} : A/J \oplus J/J^{2} \longrightarrow A/J^{2}\) 表示映射 \((x,y) \mapsto h(x)+y\)。

命题 2. - 赋予 \(k\)-模 \(A / J \oplus J / J^2\) 以上面例子中定义的 \(k\)-代数结构。映射 \(v \mapsto h_v\) 与 \(h \mapsto \psi_h\) 诱导出下列集合之间的双射：

\begin{enumerate}
\def\labelenumi{(\roman{enumi})}
\item
  \(\bar{d}\) 的 A/J-线性收缩 v 组成的集合；
\item
  k-代数同态 \(h: A/J \to A/J^{2}\) 中满足 \(\pi \circ h = Id_{A/J}\) 者组成的集合；
\item
  k-代数同构 \(\psi : \mathrm{A} / \mathrm{J} \oplus \mathrm{J} / \mathrm{J}^2 \longrightarrow \mathrm{A} / \mathrm{J}^2\)（满足 \(\pi \circ \psi = \mathrm{pr}_1\) 且 \(\psi(0, z) = z\)，对 \(z \in \mathrm{J} / \mathrm{J}^2\)）组成的集合。
\end{enumerate}

以 \(C = A / J^2\) 与 \(N = J / J^2\) 应用命题 1。设 \(\varphi : A \to A / J\) 为典范满射；同态 \(\rho\) 是 \(\varphi\) 到 \(A / J^2\) 中的一个提升。A-模 \(\text{Hom}_{A / J}((A / J) \otimes_A \Omega_k(A), J / J^2)\) 与 \(\text{Hom}_A(\Omega_k(A), J / J^2)\) 恒同；由命题 1，映射 \(v \mapsto H_v\) 是从这个集合到 \(\varphi\) 到 \(A / J^2\) 中的提升集合上的双射。为证 \(x \in J\) 时有 \(1 \otimes dx = \bar{d}(\rho(x))\)；为使 \(H_v\) 在 \(J\) 上为零，必须且只需 \(v \circ \bar{d}\) 是 \(J / J^2\) 的恒同映射。这就证明了映射 \(v \mapsto h_v\) 在陈述中所述前两个集合之间诱导出双射。

映射 \(h \mapsto \psi_h\) 是从 \(k\)-线性映射（自 A/J 到 A/J \(^2\)）组成的集合，到 \(k\)-线性映射 \(\psi : \text{A/J} \oplus \text{J/J}^2 \longrightarrow \text{A/J}^2\)（满足 \(\psi(0,z) = z\)，对 \(z \in \text{J/J}^2\)）组成的集合上的双射；此外，为使 \(\pi \circ \psi_h = \text{pr}_1\)，必须且只需 \(\pi \circ h = \text{Id}_{\text{A/J}}\)，即有 \(z \equiv h(\pi(z)) \pmod{\text{J/J}^2}\)（对每个 \(z \in \text{A/J}^2\)）。设这些条件满足。为使 \(h\) 是环同态，必须且只需 \(\psi_h\) 亦然；此外，同态 \(\psi_h\) 是双射的：其逆映射把元素 \(z\)（A/J \(^2\) 中的）对应为偶 \((\pi(z), z - h(\pi(z)))\)。这就证明了映射 \(h \mapsto \psi_h\) 在陈述中所述后两个集合之间诱导出双射。

